\documentclass[12pt,a4paper]{article}

\usepackage[T1]{fontenc}
\usepackage[utf8]{inputenc}
\usepackage[brazil]{babel}
\usepackage{lmodern}
\usepackage{geometry}
\usepackage{hyperref}
\usepackage{enumitem}
\usepackage{array}

\geometry{margin=2.5cm}
\hypersetup{
  colorlinks=true,
  linkcolor=blue,
  urlcolor=blue
}

\title{Plano de Aula: Aula 2 -- Processo de Desenvolvimento de Software}
\date{}

\begin{document}
\maketitle

\noindent
\textbf{Disciplina:} Engenharia de Software (Curso Técnico em Desenvolvimento de Sistemas)\\
\textbf{Duração:} 2 horas (120 minutos)\\
\textbf{Modalidade:} Presencial\\
\textbf{Materiais da aula:} \href{aula-02.html}{slides (HTML)} \quad \href{aula-02.pdf}{ementa (PDF)}\\
\textbf{Livro-base (referência):} \href{../99-Materiais/Engenharia%20de%20software%20projetos%20e%20processos%20by%20Wilson%20de%20P%C3%A1dua%20Paula%20Filho%20.epub.pdf}{Engenharia de software: projetos e processos (PDF)}

\section{Competências e Habilidades}
\textbf{Competências}
\begin{itemize}[leftmargin=*,itemsep=0.25em]
  \item Aplicar princípios da engenharia de software para desenvolvimento de sistemas.
  \item Compreender processos e modelos de desenvolvimento para organizar entregas.
\end{itemize}

\textbf{Habilidades}
\begin{itemize}[leftmargin=*,itemsep=0.25em]
  \item Compreender metodologias de desenvolvimento de software.
  \item Identificar atividades, papéis, artefatos e critérios de um processo.
  \item Comparar modelos de ciclo de vida e escolher o mais adequado ao cenário.
\end{itemize}

\section{Objetivos da Aula}
\begin{itemize}[leftmargin=*,itemsep=0.25em]
  \item Entender o que é um processo de software e por que ele existe.
  \item Diferenciar processo, método e ferramenta.
  \item Conhecer e comparar modelos de ciclo de vida (cascata, incremental, iterativo, espiral).
  \item Identificar riscos e consequências de processos inadequados ou inexistentes.
\end{itemize}

\section{Assuntos Abordados no Dia}
\begin{itemize}[leftmargin=*,itemsep=0.25em]
  \item Conceitos: processo, método e ferramenta; papéis, atividades, artefatos e critérios.
  \item Ciclo de vida do software e modelos de processo.
  \item Como escolher um modelo: escopo, risco, mudança, equipe e domínio.
  \item Atividade: comparar modelos e discutir vantagens, limitações e riscos.
\end{itemize}

\section{Metodologia}
Exposição dialogada com exemplos e diagramas (cascata vs.\ iterativo/incremental). Discussão em pequenos grupos e plenária curta para comparação de modelos. O professor conduz perguntas que forçam justificativa (por que um modelo é melhor em determinado cenário) e conecta com riscos e controle de mudanças.

\section{Materiais e Recursos}
\begin{itemize}[leftmargin=*,itemsep=0.25em]
  \item Quadro e marcadores.
  \item Slides: \href{aula-02.html}{aula-02.html} e ementa em PDF: \href{aula-02.pdf}{aula-02.pdf}.
  \item Livro-base (PDF): \href{../99-Materiais/Engenharia%20de%20software%20projetos%20e%20processos%20by%20Wilson%20de%20P%C3%A1dua%20Paula%20Filho%20.epub.pdf}{Engenharia de software: projetos e processos}.
  \item Handout para grupos (opcional): 4 cenários de projeto para escolha de processo.
  \item Vídeos complementares (links nos slides).
\end{itemize}

\section{Cronograma e Descrição Detalhada (120 Minutos)}
\subsection*{Parte 1: Abertura e objetivos (10 min)}
\begin{itemize}[leftmargin=*,itemsep=0.35em]
  \item Retomar a Aula 1: por que planejamento importa.
  \item Apresentar os objetivos do encontro e a atividade de comparação.
\end{itemize}

\subsection*{Parte 2: Processo, método e ferramenta (25 min)}
\begin{itemize}[leftmargin=*,itemsep=0.35em]
  \item Definir \textbf{processo}: fluxo de trabalho, papéis, artefatos e critérios.
  \item Definir \textbf{método}: técnicas e práticas (ex.: Scrum, UML).
  \item Definir \textbf{ferramenta}: suporte ao método/processo (tracker, CI, repositório).
  \item Exemplos de falhas por confusão: ``ritual sem processo'' e ``ferramenta sem método''.
\end{itemize}

\subsection*{Parte 3: Modelos de ciclo de vida (35 min)}
\begin{itemize}[leftmargin=*,itemsep=0.35em]
  \item Apresentar cascata, incremental, iterativo e espiral (objetivos e características).
  \item Discutir riscos: mudança de requisito, integração tardia, feedback tardio.
  \item Critérios de escolha: risco, incerteza, complexidade, requisitos, equipe e ambiente.
\end{itemize}

\subsection*{Parte 4: Atividade em grupos (30 min)}
\begin{itemize}[leftmargin=*,itemsep=0.35em]
  \item Dividir em grupos (3--4).
  \item Cada grupo escolhe um cenário (ex.: sistema escolar, e-commerce pequeno, software crítico).
  \item Responder:
  \begin{itemize}[leftmargin=*,itemsep=0.25em]
    \item Qual modelo de processo escolher e por quê.
    \item Quais riscos existem e como o modelo ajuda (ou atrapalha).
    \item Quais artefatos mínimos produzir (requisitos, plano, testes).
  \end{itemize}
\end{itemize}

\subsection*{Parte 5: Socialização e fechamento (20 min)}
\begin{itemize}[leftmargin=*,itemsep=0.35em]
  \item 2--3 grupos apresentam (2--3 min cada).
  \item Professor consolida: não há modelo perfeito; há adequação ao contexto.
  \item Ponte para a Aula 3: como métodos ágeis se encaixam em processos iterativos/incrementais.
\end{itemize}

\section{Avaliação}
\begin{itemize}[leftmargin=*,itemsep=0.25em]
  \item Participação na discussão e justificativa da escolha do processo.
  \item Qualidade do raciocínio sobre riscos, artefatos mínimos e critérios.
\end{itemize}

\section{Referências}
\begin{itemize}[leftmargin=*,itemsep=0.25em]
  \item \href{../99-Materiais/Engenharia%20de%20software%20projetos%20e%20processos%20by%20Wilson%20de%20P%C3%A1dua%20Paula%20Filho%20.epub.pdf}{PAULA FILHO, Wilson de Pádua. Engenharia de software: projetos e processos. 4. ed. 2019.}
\end{itemize}

\end{document}
